\subsection{复合扩张的应用}

命题 10. -- 设 L 和 M 是域 K 的两个扩张，(E, u, v) 是 L 和 M 的一个复合扩张（V，第 12 页）。假设环 L ⊗\_K M 是整环，且 E 的子扩张 u(L) 和 v(M) 在 K 上代数不相交。则 u(L) 和 v(M) 在 K 上线性不相交。

令 \(w = u * v\) （V，第 12 页），用 F 表示整环 \(L \otimes_{K} M\) 的分式域，并借助映射 \(x \mapsto x \otimes 1\) （相应地， \(y \mapsto 1 \otimes y\) ）把 L（相应地，M）等同于 F 的一个子域；则 w 在 L（相应地，M）上的限制是 u（相应地，v）。设 B 是 M 在 K 上的一个超越基（V，第 109 页，定理 1）。

根据假设， \(u(\mathbf{L})\) 和 \(v(\mathbf{M})\) 在 \(\mathbf{K}\) 上代数不相交；因此（V，第 114 页，命题 14）， \(u(\mathbf{L})\) 和 \(v(\mathbf{K}(\mathbf{B}))\) 在 \(\mathbf{K}\) 上线性不相交。于是存在一个 K-同态 \(u': \mathbf{L}(\mathbf{B}) \to \mathbf{E}\) ，它与 \(u\) 在 \(\mathbf{L}\) 上一致，且与 \(v\) 在 \(\mathbf{K}(\mathbf{B})\) 上一致。根据构造， \(\mathbf{L}\) 和 \(\mathbf{M}\) 在 \(\mathbf{K}\) 上于 \(\mathbf{F}\) 中线性不相交；由命题 8（V，第 15 页），子域 \(\mathbf{L}(\mathbf{B})\) 和 \(\mathbf{M}\) （\(\mathbf{F}\) 的）在 \(\mathbf{K}(\mathbf{B})\) 上线性不相交。由此可知，存在一个 K-同态 \(w' \colon \mathbf{M}[\mathrm{L}(\mathrm{B})] \to \mathrm{E}\) ，它与 \(u'\) 在 \(\mathrm{L}(\mathrm{B})\) 上一致，且与 \(v\) 在 \(\mathbf{M}\) 上一致。但域 \(F\) 由 \(\mathbf{M} \cup \mathrm{L}(\mathrm{B})\) 生成，且 \(\mathbf{M}\) 在 \(\mathbf{K}(\mathbf{B})\) 上是代数的；于是我们有 \(\mathbf{M}[\mathrm{L}(\mathbf{B})] = F(\mathbf{V}, p. 18, Cor. 2)\) 。由此我们得出结论： \(w'\) 是 \(F\) 到 \(E\) 上的一个 K-同构，它在 \(L\) （相应地， \(M\) ）上的限制是 \(u\) （相应地， \(v\) ）。这表明 \(u(L)\) 和 \(v(M)\) 在 \(K\) 上线性不相交。

推论 1. --- 设 \(\Omega\) 是域 \(K\) 的一个扩张， \(L\) 是 \(\Omega\) 的一个在 \(K\) 上正则的子扩张。每个子扩张 \(M\) （\(\Omega\) 的），只要它与 \(L\) 在 \(K\) 上代数不相交，就是线性不相交的。

由正则扩张的定义，环 \(L \otimes_{K} M\) 是一个整环，现在只需应用命题 10。

推论 2. --- 设 \(\Omega\) 是域 \(K\) 的一个扩张， \(L\) , \(M\) 是 \(\Omega\) 的两个子扩张。假设 \(L\) 在 \(K\) 上可分，且 \(K\) 在 \(M\) 中的相对可分闭包等于 \(K\) 。若 \(L\) 和 \(M\) 在 \(K\) 上代数不相交，则它们在 \(K\) 上线性不相交。

根据命题 10，只需注意到环 \(\mathbf{L} \otimes_{\mathbf{K}} \mathbf{M}\) 是整环（V，第 140 页，推论）。

